% =========================================================================
\chapter{Direct \& Factorization-Based Silicon Matrix Inverters}
\label{ch:direct_solvers}
% =========================================================================

\section{Introduction to Matrix Factorization Hardware}
While block tile solvers (like MI2D) accelerate dense square matrices, hardware chips for physical AI, 5G/6G wireless communication, and embedded signal processing frequently require **direct factorization** of structured, symmetric positive-definite, or ill-conditioned matrices.

Matrix factorization decomposes a complex matrix $A$ into product factors of simpler matrices (such as lower/upper triangular, orthogonal, or diagonal matrices):
\begin{itemize}[leftmargin=*]
    \item \textbf{QR Decomposition}: $A = Q R$ (where $Q$ is orthogonal $Q^T Q = I$, and $R$ is upper triangular).
    \item \textbf{Cholesky / $L D L^T$ Decomposition}: $A = L L^T$ or $A = L D L^T$ (for symmetric positive-definite matrices).
    \item \textbf{State-Space Realization}: Factorizing time-varying structured operators into state updates.
    \item \textbf{Iterative Solvers}: Jacobi and Gauss-Seidel iterations on FPGA logic.
\end{itemize}

% -------------------------------------------------------------------------
\section{QR Decomposition via Givens Rotations \& CORDIC Pipelines}
% -------------------------------------------------------------------------
QR decomposition factors matrix $A \in \mathbb{R}^{M \times N}$ into an orthogonal matrix $Q \in \mathbb{R}^{M \times M}$ and an upper triangular matrix $R \in \mathbb{R}^{M \times N}$.

To solve linear systems $A x = b$:
\begin{equation}
Q R x = b \implies R x = Q^T b
\end{equation}
Since $R$ is upper triangular, $x$ is evaluated directly via **back-substitution** without matrix inversion.

\subsection{Givens Rotations}
A **Givens Rotation Matrix** $G(i, j, \theta)$ rotates vectors in the $(i, j)$ plane to zero out a target off-diagonal matrix element:

\begin{equation}
G(i, j, \theta) = \begin{bmatrix}
1 & \dots & 0 & \dots & 0 & \dots & 0 \\
\vdots & \ddots & \vdots & & \vdots & & \vdots \\
0 & \dots & c & \dots & s & \dots & 0 \\
\vdots & & \vdots & \ddots & \vdots & & \vdots \\
0 & \dots & -s & \dots & c & \dots & 0 \\
\vdots & & \vdots & & \vdots & \ddots & \vdots \\
0 & \dots & 0 & \dots & 0 & \dots & 1
\end{bmatrix}
\end{equation}
where $c = \cos\theta = \frac{x_i}{\sqrt{x_i^2 + x_j^2}}$ and $s = \sin\theta = \frac{x_j}{\sqrt{x_i^2 + x_j^2}}$.

Applying a sequence of Givens rotations systematically annihilates lower triangular elements, producing $R = G_k \dots G_2 G_1 A$.

\subsection{CORDIC Algorithm Hardware Pipeline}
Computing trigonometric functions ($\cos\theta, \sin\theta$) and square roots ($\sqrt{x_i^2 + x_j^2}$) directly in digital logic consumes extensive silicon area. 

FPGA hardware accelerators replace floating-point division and square root units with the **CORDIC (Coordinate Rotation Digital Computer)** algorithm. CORDIC evaluates rotations using only **shift-and-add operations**:

\begin{align}
x_{k+1} &= x_k - d_k \cdot y_k \cdot 2^{-k} \\
y_{k+1} &= y_k + d_k \cdot x_k \cdot 2^{-k} \\
z_{k+1} &= z_k - d_k \cdot \arctan(2^{-k})
\end{align}
where $d_k \in \{-1, +1\}$ is the direction of rotation. 

A fully pipelined CORDIC array on FPGA computes Givens QR rotations in **one clock cycle per iteration step**, offering high numerical stability for ill-conditioned matrices.


% -------------------------------------------------------------------------
\section{Cholesky ($L L^T$) \& $L D L^T$ Decomposition Engines for MIMO Systems}
% -------------------------------------------------------------------------
In 5G wireless multi-antenna (MIMO) detection and radar signal processing, Minimum Mean Square Error (MMSE) detectors solve:

\begin{equation}
x_{\text{MMSE}} = (H^H H + \sigma^2 I)^{-1} H^H y
\end{equation}
The matrix $A = (H^H H + \sigma^2 I)$ is **Hermitian positive-definite**.

\subsection{$L D L^T$ Factorization}
Lingala et al. (IISc Bangalore 2022)~\cite{lingala2022matrix} developed an FPGA hardware architecture for matrix inversion using $L D L^T$ decomposition:

\begin{equation}
A = L D L^T
\end{equation}
where $L$ is a unit lower triangular matrix (ones on the diagonal), and $D$ is a purely diagonal matrix.

Inverting $A$ using $L D L^T$:
\begin{equation}
A^{-1} = (L^T)^{-1} D^{-1} L^{-1}
\end{equation}

Why $L D L^T$ is preferred on FPGAs:
\begin{itemize}[leftmargin=*]
    \item **Eliminates Square Root Units**: Standard Cholesky ($A = L L^T$) requires evaluating square roots $\sqrt{L_{ii}}$ on diagonal elements. $L D L^T$ eliminates square root operations entirely, using simple scalar divisions $\frac{1}{D_{ii}}$.
    \item **Low Compute Load**: Lingala et al. implemented $L D L^T$ using iterative $2 \times 2$ block inversions, achieving low computational overhead on Xilinx Zynq FPGA boards.
\end{itemize}

\subsection{SISO MMSE MIMO ASIC Inverters}
Auras et al. (IEEE ISCAS 2014)~\cite{auras2014efficient} proposed a dedicated Application-Specific Integrated Circuit (ASIC) for matrix inversion in soft-input soft-output MMSE MIMO detectors. 

Auras et al. achieved the IEEE 802.11n peak data rate of **$600\,\text{Mbit/s}$**, outperforming competing VLSI architectures by a factor of **$1.7\times$** in throughput-per-area efficiency.


% -------------------------------------------------------------------------
\section{State-Space Realization for Structured Matrices}
% -------------------------------------------------------------------------
As established by van der Veen and Dewilde (1993)~\cite{van1993large}, large structured matrices originating from time-series signals or finite-element physical grids can be represented as the input-output operator of a **time-varying state-space system**:

\begin{align}
x_{k+1} &= A_k x_k + B_k u_k \\
y_k &= C_k x_k + D_k u_k
\end{align}

Van der Veen and Dewilde demonstrated that if the number of internal states is small relative to matrix dimension $N$, the structured matrix inverse $A^{-1}$ can be computed in $\mathcal{O}(N \cdot s^2)$ time (where $s \ll N$ is the state dimension) using continuous state realizations, bypassing standard $\mathcal{O}(N^3)$ dense inversion entirely!


% -------------------------------------------------------------------------
\section{Iterative Solvers: FPGA Jacobi Load Flow Engines}
% -------------------------------------------------------------------------
For large sparse matrices (such as electrical power grid networks), direct matrix inversion introduces fill-in zero redundancy. 

Foertsch et al. (2005)~\cite{foertsch2005jacobi} re-examined the **Jacobi Iterative Method** for fully pipelined FPGA hardware implementation:

\begin{equation}
x^{(k+1)} = D^{-1} \left( b - (L + U) x^{(k)} \right)
\end{equation}

Because $D^{-1}$ is purely diagonal, computing $D^{-1}$ requires simple scalar multiplications. Foertsch et al. built a fully pipelined FPGA hardware engine that executes Jacobi load flow updates in parallel, demonstrating that FPGA-accelerated Jacobi solvers match or exceed CPU Newton-Raphson software speeds for large power grids.


% -------------------------------------------------------------------------
\section{Comprehensive Hardware Strategy Comparison}
% -------------------------------------------------------------------------
Table~\ref{tab:ch08_hardware_comparison} presents a theoretical comparison of silicon matrix inversion strategies.

\begin{table}[htbp]
\centering
\caption{Comparison of Silicon Hardware Matrix Inversion Strategies}
\label{tab:ch08_hardware_comparison}
\begin{tabularx}{\textwidth}{>{\raggedright\arraybackslash}p{3.0cm} >{\raggedright\arraybackslash}p{2.0cm} >{\raggedright\arraybackslash}p{2.2cm} >{\raggedright\arraybackslash}X >{\raggedright\arraybackslash}p{3.0cm}}
\toprule
\textbf{Algorithm} & \textbf{Precision} & \textbf{FLOP Latency} & \textbf{FPGA Hardware Resources} & \textbf{Target Physical Application} \\
\midrule
Direct LU / Cholesky & Exact & $\mathcal{O}(N^3)$ & High (DSP Slices \& BRAM) & High-precision scientific computing \\
Givens QR (CORDIC) & High (Stable) & $\mathcal{O}(N^2)$ & Moderate (CORDIC Shift-Adders) & 5G/6G MIMO Signal Detection \\
$L D L^T$ Decomposition & High & $\mathcal{O}(N^3/6)$ & Low-to-Moderate (No Sqrt units) & Embedded MIMO Board Engines \\
MI2D (Tile Schur) & High & $\mathcal{O}(N^3/k^2)$ & Parallel PE Tile Grid Arrays & Deep Learning Newton Optimizers \\
State-Space Realization & Exact/Approx & $\mathcal{O}(N s^2)$ & Low Logic LUTs & Time-Varying Adaptive Control \\
Jacobi Iterative & Approximate & Iterative & Very Low Logic LUTs & Power Flow \& Grid Simulations \\
\bottomrule
\end{tabularx}
\end{table}
